% ECONOMICS_PROBLEM: {"id": "C24-509", "title": "Nonrandom attrition and missing-not-at-random outcomes", "jel_code": "C24", "source_id": "OP-0109"}
% FIRST_POSTED: 2026-10-09
% ATTEMPT_STATUS: partial
% ENTRY_KIND: research_attempt
\documentclass[11pt,a4paper]{article}
\usepackage[T1]{fontenc}
\usepackage[utf8]{inputenc}
\usepackage[margin=27mm]{geometry}
\usepackage{amsmath,amssymb,amsthm,mathtools}
\usepackage[hidelinks]{hyperref}
\setlength{\parindent}{0pt}
\setlength{\parskip}{5pt}
\newtheorem{theorem}{Theorem}[section]
\newtheorem{proposition}[theorem]{Proposition}
\newtheorem{lemma}[theorem]{Lemma}
\newcommand{\R}{\mathbb R}
\newcommand{\E}{\mathbb E}
\DeclareMathOperator{\Var}{Var}
\DeclareMathOperator{\Cov}{Cov}
\begin{document}
\section*{Sharp attrition bounds with bounded monotonicity violations and uniform inference}
\section{Precise claim}
This attempt completely derives the no-covariate specialization of the catalogue's bounded-outcome, independent-trial target, and gives an exact finite-covariate extension. It supplies sharp population and always-observed bounds for every defier allowance $\delta$, including mass points, and a finite-dimensional confidence construction valid uniformly as observation probabilities approach zero. The unrestricted covariate-sensitive catalogue target remains partial. It does not identify missing-not-at-random outcomes without restrictions, claim optimal confidence width, or assert novelty over the selection-bounds literature.

Assume i.i.d. units, unconditional random assignment $D$ with positive known propensity, consistency, and potential outcomes $Y_0,Y_1\in[a,b]$, $a<b$. The potential indicators $R_0,R_1$ may depend arbitrarily on both outcomes. The sensitivity class imposes only
$\Pr(R_0=1,R_1=0)\le\delta$. Through the first confidence theorem, assume the recorded baseline covariate $X$ is constant. Equivalently these formulas use the coarsened observable law that omits $X$; they need not be sharp relative to a richer law retaining informative baseline covariates. A later section explicitly incorporates finitely many covariate values.

Let $H_d$ be the observed subprobability distribution of $Y_d$ on $R_d=1$, of mass $p_d$, and $M_d=\int y\,dH_d(y)$. These are identified by randomization, even when the conditional mean given observation is undefined. Write $q=\Pr(R_0=R_1=1)$ and
\[
 q_- =\max\{0,p_0+p_1-1,p_0-\delta\},\qquad
 q_+=\min(p_0,p_1).                                      \tag{1}
\]
The four observation-stratum probabilities are $q,p_0-q,p_1-q,1-p_0-p_1+q$. Their nonnegativity and the defier bound prove that compatibility is equivalent to $q_-\le q_+$. A principal-stratum effect is defined only when $q>0$; if $q_+=0$, report it as undefined rather than zero.

\section{Sharp population bounds}
\begin{theorem}
Whenever the sensitivity model is compatible, the sharp population ATE interval is
\[
 \left[M_1+(1-p_1)a-M_0-(1-p_0)b,\quad
 M_1+(1-p_1)b-M_0-(1-p_0)a\right].         \tag{2}
\]
It does not shrink with $\delta$ alone.
\end{theorem}
\begin{proof}
Each unobserved potential-outcome contribution ranges between its missing probability times $a$ and times $b$. This proves necessity. Choose any feasible $q$. Split $H_d$ arbitrarily into an always-observed submeasure of mass $q$ and its remaining observed stratum. Couple the two always-observed outcome margins arbitrarily; in other strata assign all missing potential outcomes to the endpoints needed for the desired bound. This preserves the full observed law and all observation-stratum probabilities and attains each endpoint. Mixing the two joint latent laws preserves the observable law and the defier inequality and attains every intermediate ATE.
\end{proof}
The restriction on observation monotonicity changes who could be always observed, but imposes no value restriction on a missing outcome. This is why the familiar trimming improvement concerns a different population.

\section{Sharp principal-stratum bounds, including atoms}
When $p_d>0$, let $Q_d(v)=\inf\{y:H_d(({-\infty},y])/p_d\ge v\}$ be the generalized quantile. For $0<q\le p_d$, define
\[
 \ell_d(q)=\frac{p_d}{q}\int_0^{q/p_d}Q_d(v)\,dv,
 \qquad
 u_d(q)=\frac{p_d}{q}\int_{1-q/p_d}^1Q_d(v)\,dv.             \tag{3}
\]
Quantile integrals automatically permit fractional allocation of mass at a trimming threshold.
\begin{lemma}[Submeasure trimming]
Among submeasures $K\le H_d$ of mass $q$, the smallest and largest conditional means $q^{-1}\int y\,dK$ are $\ell_d(q)$ and $u_d(q)$, respectively. Both are attained.
\end{lemma}
\begin{proof}
For the minimum choose all mass below a $q$-quantile cutoff, the required fraction of an atom at the cutoff, and no mass above. Compare any competitor to this measure: both have mass $q$, and transferring mass from below the cutoff to above cannot decrease the integral. Formally subtract the cutoff times the zero total-mass difference; every remaining contribution is nonnegative. Reversing the ordering proves the maximum. Bounded support ensures the integrals exist. The argument applies to arbitrary distributions, not only to densities.
\end{proof}
\begin{theorem}
For fixed feasible $q>0$, the sharp interval for
$\theta_A=E[Y_1-Y_0\mid R_0=R_1=1]$ is
\[
 [L(q),U(q)]=[\ell_1(q)-u_0(q),\quad u_1(q)-\ell_0(q)].      \tag{4}
\]
If $q_->0$, the overall sharp interval is $[L(q_-),U(q_-)]$. If $q_-=0<q_+$, its exact set is
$\bigcup_{0<q\le q_+}[L(q),U(q)]$, with closure endpoints given by the limits as $q\downarrow0$.
\end{theorem}
\begin{proof}
The always-observed outcome margins must be submeasures of $H_0,H_1$ of mass $q$, so the lemma gives necessity. To attain the lower endpoint, select the bottom treated submeasure and top control submeasure and couple their normalized distributions in the always-observed stratum. Put the leftover observed mass in the treated-only and control-only strata. Fill missing potential outcomes arbitrarily and use the stratum probabilities from (1). This is a compatible latent law. Reverse top and bottom for the upper endpoint. Mixing these two laws keeps $q$ fixed and attains the entire interval.

The average of the bottom $q$ mass is nondecreasing in $q$, and the average of the top $q$ mass is nonincreasing. This follows directly from the ordered quantile integrals: newly added bottom-tail values are at least its previous average, and newly added top-tail values are at most its previous average. Consequently $L(q)$ is nondecreasing and $U(q)$ nonincreasing. All intervals are nested. Their largest member occurs at $q_-$ if that value is positive; otherwise their union has the stated limiting endpoints. At $q_-=0$, an endpoint may be only approachable, not attained.
\end{proof}
The limiting lower endpoint is $\operatorname{ess\,inf}(Y_1\mid R_1=1)-\operatorname{ess\,sup}(Y_0\mid R_0=1)$, and the upper endpoint reverses extrema. This statement is used only when both observed distributions have positive mass. With $\delta=0$ and compatible $p_1\ge p_0>0$, (1) forces $q=p_0$, yielding the usual monotone-selection trimming formula.

As a numerical check, let $p_0=0.6,p_1=0.8$, with observed binary outcomes having conditional success probabilities $0.5$ and $0.75$. At $\delta=0$, $q=0.6$, the control mean is $0.5$ and treated trimmed means range from $2/3$ to one, yielding $[1/6,1/2]$. At $\delta=0.2$, $q_-=0.4$; control trimmed means range from $0.25$ to $0.75$, and treated means from $0.5$ to one, yielding $[-0.25,0.75]$. Formula (2) gives population bounds $[-0.1,0.5]$ for either allowance. These intervals answer distinct targets.

\section{Uniform, simultaneous, computable confidence regions}
Scale outcomes to $Z=(Y-a)/(b-a)\in[0,1]$. Choose a deterministic integer $K\ge1$ before seeing outcomes and round down to $Z^K=\lfloor KZ\rfloor/K$, with $Z^K=1$ when $Z=1$. The arm's observable category is either missing or one of the $K+1$ observed grid values. Let $\widehat v_{dc}$ be its empirical category frequency among the $n_d$ assigned units. For $n_d>0$ set
\[
 \epsilon_d=\sqrt{\log(4(K+2)/\alpha)/(2n_d)},\qquad
 I_{dc}=[\widehat v_{dc}-\epsilon_d,\widehat v_{dc}+\epsilon_d]\cap[0,1]. \tag{5}
\]
For $n_d=0$ set every category interval to $[0,1]$. Conditional on assignment counts, Hoeffding and a union bound over $2(K+2)$ categories show that all true category probabilities lie in their intervals with probability at least $1-\alpha$. This conditional assertion also holds unconditionally. No lower bound on $p_d$ or density condition is used.

Introduce $4(K+1)^2$ nonnegative latent cell masses
$x_{r_0r_1jk}$, for the two observation indicators and grid potential outcomes $j/K,k/K$. Require their sum to be one. For each arm, sum the cells producing each observable category and constrain that sum to the appropriate $I_{dc}$. Impose
\[
 \sum_{j,k}x_{10jk}\le\delta,\qquad
 q(x)=\sum_{j,k}x_{11jk}.                                 \tag{6}
\]
Call the resulting finite polytope $\mathcal X_\delta$. Define
\[
 A(x)=\sum_{r_0,r_1,j,k}x_{r_0r_1jk}(k-j)/K,
 \quad N(x)=\sum_{j,k}x_{11jk}(k-j)/K.                       \tag{7}
\]
The population confidence interval is the range of $A(x)$ over $\mathcal X_\delta$, enlarged by $2/K$ in each direction and intersected with $[-1,1]$. For the principal effect, range $N(x)/q(x)$ over feasible $q(x)>0$, enlarge by $2/K$, and intersect with $[-1,1]$. Multiply all effect coordinates by $b-a$ to return to original units.

A joint region retains pairs within $2/K$ in each coordinate of
$(A(x),N(x)/q(x))$ for one common feasible $x$ with $q(x)>0$. This preserves dependence between the two targets; the Cartesian product of the separate intervals is a simpler conservative joint region. If feasible $q=0$ occurs, separately report that existence of the principal stratum is not certified. If no feasible positive $q$ exists, the principal effect is undefined in the retained class. An empty polytope reports rejection at that sensitivity value, not a numerical zero effect.

\begin{theorem}[Finite-sample uniform coverage]
With probability at least $1-\alpha$, simultaneously for every $\delta\in[0,1]$, the regions just described contain the entire identified target set of every latent model compatible with the true observable law and that defier allowance. Principal coverage applies only to models with $q>0$. The assertion is uniform over the stated i.i.d. bounded-outcome model, including mass points and selection rates tending to zero.
\end{theorem}
\begin{proof}
On the single category event from (5), round both potential outcomes in any compatible latent model. Its resulting cell masses satisfy every observable interval and (6), hence belong to $\mathcal X_\delta$. Rounding preserves the observation indicators, the defier probability and $q$. Each outcome changes by at most $1/K$, so both its population difference and its conditional difference given the always-observed stratum change by at most $2/K$. This bound has no factor $1/q$: it bounds individual differences before taking the conditional mean. Thus each true target lies in the stated enlargement. The same event works for every compatible model and every $\delta$; no union bound over sensitivity values is required.
\end{proof}

The population extrema use linear programming. Principal extrema are linear-fractional programs with a finite linear reformulation: set $y=x/q$, $t=1/q$, impose $\sum y=t$, $\sum y_{11jk}=1$, and multiply each probability constraint's constant by $t$. The objective is the linear numerator in $y$. This computes sharp infima and suprema for the retained grid models. For a fixed candidate joint effect pair, multiply its principal-effect inequalities by $q(x)$ and impose the resulting closed linear inequalities along with the population-effect inequalities. Maximize $q(x)$ over this compact feasible polytope and accept membership exactly when its optimum is positive. This explicitly handles the strict $q>0$ condition using an ordinary linear program. Thus the confidence construction is operational, rather than merely a definition over unspecified models.

\section{Informativeness and unavoidable weak-selection limits}
A fixed coarse grid is honest but needlessly conservative. For total sample size $n$, choose for example $K_n=\max(1,\lfloor n^{1/4}\rfloor)$. With treatment propensity bounded away from zero and one, the grid error vanishes and the total category-mass uncertainty is
$O_P(K_n\sqrt{\log K_n/n})=o_P(1)$. At any compatible law with $q_->0$, the trimmed endpoints consequently converge to their population values. To see why atoms cause no difficulty, the lower unnormalized trimmed integral on $[0,1]$ can be written
\[
 L_H(q)=\max_{t\in[0,1]}\{qt-\int(t-z)_+\,dH(z)\}.
\]
The integral equals $\int_0^tH([0,z])\,dz$, so perturbing the sub-CDF uniformly and perturbing $q$ changes $L_H(q)$ by at most the sum of those perturbations. The upper integral has the analogous minimum representation. Dividing by $q$ bounded away from zero preserves continuity. The category and grid bounds control these perturbations. This establishes nonvacuous consistency in the regular positive-stratum regime, without claiming an optimal convergence rate.

Uniformly shrinking principal intervals near $q=0$ are impossible in general. Let $R_0=R_1$ have probability $q_n$ of being one, $Y_0=0$, and compare two models with $Y_1=0$ versus $Y_1=1$ on that rare stratum. Their principal effects differ by one, but their full samples can be coupled to agree unless at least one unit from that stratum is observed, an event of probability at most $nq_n$. If $nq_n\to0$, no data-based procedure can reliably separate the two. A uniformly honest region must remain wide with substantial probability. The construction above accepts this limitation instead of dividing estimated trimmed sums by a nearly zero estimated rate.

An independent numerical check compared both population and both principal extrema against direct latent-cell linear programs in 400 pseudorandom three-point outcome examples, including varying defier allowances and zero lower stratum bounds. All 1,600 extrema agreed with the analytic formulas to numerical tolerance. The principal programs used the linear-fractional transformation above; this is a finite validation of the proofs, not a substitute for them.

\section{Observed covariates change the sharp set}
Ignoring baseline covariates can discard restrictions even under unconditional assignment. Suppose half the population has $X=0$ and is always observed in both regimes, while the other half has $X=1$ and is never observed. The full observed law forces $q=1/2$. At $\delta=1/2$, the aggregate observation probabilities alone instead allow any $q\in[0,1/2]$. Trimming the aggregate outcome distributions can therefore give a strictly wider principal-effect interval. No assumption that covariates are merely recorded removes their common cross-regime identity.

For finitely many baseline groups $g$ with probabilities $w_g>0$, let $p_{dg}$ and the trimmed functions $\ell_{dg},u_{dg}$ be defined conditionally within group. Choose conditional always-observed masses
\[
 \max(0,p_{0g}+p_{1g}-1)\le q_g\le\min(p_{0g},p_{1g}),
 \quad \sum_gw_g(p_{0g}-q_g)\le\delta,
 \quad Q=\sum_gw_gq_g>0.
\]
The sharp lower principal endpoint is the infimum over these masses of
\[
 \frac{\sum_gw_gq_g[\ell_{1g}(q_g)-u_{0g}(q_g)]}{Q},
\]
and the upper endpoint is the corresponding supremum with $u_{1g}-\ell_{0g}$. A group with $q_g=0$ contributes zero and requires no quantile evaluation. Necessity follows from conditional trimming. Sufficiency follows by constructing the four observation strata and their extremal outcome submeasures separately in each group, then mixing groups with weights $w_g$. Every retained latent model reproduces the full joint law with $X$, and intermediate effects are attained by mixtures preserving that law. Unlike the constant-$X$ case, increasing $Q$ can change the composition of the forced strata, so replacing the optimization by a single aggregate minimum $q$ is generally invalid.

For a finite outcome grid, this is exactly the earlier latent-mass program with an additional group index and observable categories $(g,R,Z^K)$. With $G$ groups, replace $K+2$ by $G(K+2)$ in (5). Require one common latent group distribution across treatment arms. All cell constraints remain linear, the principal objective remains linear-fractional, and the same rounding proof establishes uniform simultaneous coverage. A fixed finite partition of continuous $X$ also gives a valid outer confidence region for the richer model, but generally loses sharpness relative to its full observable law. Efficient refinement and informative inference with arbitrary high-dimensional covariates are not established here.

\section{Resolution and scope boundary}
Equations (1)--(4) solve the constant-covariate sensitivity problem under the declared outcome support and observation-stratum restriction; the grouped extension preserves a finite observed covariate law. Equations (5)--(7) supply honest simultaneous inference with explicit computation and regular-case informativeness. Covariate-specific monotonicity, clustered dependence, undefined wages outside employment, informative treatment assignment, or extra restrictions linking principal and population effects change the model and require a separate derivation. These are complete results for the stated specializations. A sharp, practical treatment of an unrestricted covariate law and its associated inference is left open, so the overall attempt remains partial.
\begin{thebibliography}{9}
\bibitem{lee} D. S. Lee (2009), \emph{Training, Wages, and Sample Selection: Estimating Sharp Bounds on Treatment Effects}, Review of Economic Studies 76, 1071--1102. \url{https://doi.org/10.1111/j.1467-937X.2009.00536.x}. Earlier author-hosted working paper: \url{https://www.princeton.edu/~davidlee/wp/w11721.pdf}. The monotone-selection trimming benchmark is established literature; the sensitivity construction and confidence projection above are stated and proved explicitly, without a priority claim.
\end{thebibliography}

\end{document}
